\documentclass[../../main.tex]{subfiles}
\begin{document}
\section{群作用对光滑流形的商}\label{section:ch21-quotients-by-actions}

设给定群 $G$ 在拓扑空间 $M$ 上的作用。回忆点 $p\in M$ 的\emph{轨道}是 $p$ 在群所有元素作用下的像集合 $G\cdot p=\{g\cdot p\colon g\in G\}$。轨道集合记为 $M/G$，赋予商拓扑，称为该作用的\emph{轨道空间}。

\begin{lemma}\label{lemma:ch21-商映射开}
对拓扑群 $G$ 在拓扑空间 $M$ 上的任意连续作用，商映射 $\pi\colon M\to M/G$ 是开映射。
\end{lemma}

\begin{proof}
对任意 $g\in G$、任意子集 $U\subseteq M$，定义 $g\cdot U=\{g\cdot x\colon x\in U\}$。若 $U$ 开，则 $\pi^{-1}(\pi(U))$ 等于所有 $gU$（$g\in G$）之并。由于 $p\mapsto g\cdot p$ 是同胚，每个 $gU$ 开，故 $\pi^{-1}(\pi(U))$ 开。因 $\pi$ 是商映射，$\pi(U)$ 开。
\end{proof}

\begin{example}{[光滑 Lie 群作用的轨道空间]}\label{example:ch21-轨道空间}
(a) 平凡作用 $g\cdot p=p$ 的轨道是单点集，$M/G=M$。
(b) $\mathbb{R}^k$ 在 $\mathbb{R}^k\times\mathbb{R}^n$ 上平移作用 $v\cdot(x,y)=(v+x,y)$，轨道平行于 $\mathbb{R}^k$，商同胚于 $\mathbb{R}^n$。
(c) $\mathbb{S}^1$ 在 $\mathbb{C}$ 上旋转作用，轨道为同心圆及原点，商同胚于 $[0,\infty)$，原点处不是流形。
(d) $\operatorname{GL}(n,\mathbb{R})$ 在 $\mathbb{R}^n$ 上自然作用有两条轨道 $\{0\}$ 与 $\mathbb{R}^n\setminus\{0\}$，商甚至非 Hausdorff。
(e) $O(n)$ 在 $\mathbb{R}^n$ 上的作用轨道为原点与球面，商同胚于 $[0,\infty)$。
(f) 删去原点后：$(\mathbb{C}\setminus\{0\})/\mathbb{S}^1\cong\mathbb{R}^+$；$(\mathbb{R}^n\setminus\{0\})/O(n)$ 同理；$(\mathbb{R}^n\setminus\{0\})/\operatorname{GL}(n,\mathbb{R})$ 是单点。
(g) $O(n)$ 在 $\mathbb{S}^{n-1}$ 上的作用可迁，商为单点。
\end{example}

\begin{example}\label{example:ch21-无理流}
设 $\alpha$ 无理，$\mathbb{R}$ 经 $t\cdot(w,z)=(e^{2\pi\mathrm{i}t}w,e^{2\pi\mathrm{i}\alpha t}z)$ 作用于 $\mathbb{T}^2$。该作用光滑、自由、轨道稠密，故 $\mathbb{T}^2/\mathbb{R}$ 有平凡拓扑，非 Hausdorff。
\end{example}

\begin{proposition}\label{proposition:ch21-真作用Hausdorff}
若 Lie 群在流形上连续且真地作用，则轨道空间是 Hausdorff 的。
\end{proposition}

\begin{proof}
令 $\Theta(g,p)=(g\cdot p,p)$，$\pi$ 为商映射。轨道关系 $O=\Theta(G\times M)=\{(g\cdot p,p)\}$。真连续映射闭（定理 A.57），故 $O$ 闭。由引理\cref{lemma:ch21-商映射开}，$\pi$ 开，习题 A.36 给出 $M/G$ Hausdorff。
\end{proof}

\begin{proposition}{[真作用的刻画]}\label{proposition:ch21-真作用刻画}
下列各条等价：(a) 作用真；(b) 若 $(p_i)$、$(g_i)$ 使 $(p_i)$ 与 $(g_i\cdot p_i)$ 都收敛，则 $(g_i)$ 有收敛子列；(c) 对每个紧 $K\subseteq M$，$G_K=\{g\in G\colon(gK)\cap K\neq\varnothing\}$ 紧。
\end{proposition}

\begin{proof}
作用真当且仅当 $\Theta(g,p)=(g\cdot p,p)$ 真。(a)$\Rightarrow$(b)：$(g_i,p_i)$ 最终落入紧集 $\overline V\times\overline U$，故有收敛子列。(b)$\Rightarrow$(c)：若 $G_K$ 不紧，取 $(g_i)\subset G_K$ 无收敛子列，配 $p_i\in K$ 使 $g_i\cdot p_i\in K$，由 $K$ 紧与 (b) 得矛盾。(c)$\Rightarrow$(a)：若 $L\subseteq M\times M$ 紧，$K_1,K_2$ 为其投影，则 $\Theta^{-1}(L)\subseteq G_{K_1}\times K_1$ 紧且闭，故紧。
\end{proof}

\begin{corollary}\label{corollary:ch21-紧群真}
紧 Lie 群在流形上的每个连续作用都是真的。
\end{corollary}

\begin{proof}
由命题\cref{proposition:ch21-真作用刻画}(b)，$G$ 中每序列有收敛子列。
\end{proof}

\begin{proposition}\label{proposition:ch21-真作用轨道}
若 $G$ 真光滑作用于 $M$，则对 $p\in M$，轨道映射 $\sigma_p\colon G\to M$ 真，故轨道 $G\cdot p$ 闭；若 $G_p=\{e\}$，则 $\sigma_p$ 是光滑嵌入，轨道为真嵌入子流形。
\end{proposition}

\begin{proof}
$\sigma_p^{-1}(K)$ 闭且含于 $G_K\cup\{p\}$，由命题\cref{proposition:ch21-真作用刻画} 紧，故 $\sigma_p$ 真；由定理 A.57 轨道闭。嵌入结论由命题 7.26 与 4.22。
\end{proof}

\begin{corollary}\label{corollary:ch21-真作用轨道闭}
真作用下每条轨道是 $M$ 的闭子集，每个迷向群紧。
\end{corollary}

\begin{example}\label{example:ch21-相似作用}
$\mathbb{R}^+$ 经 $t\cdot(x^1,\dots,x^n)=(tx^1,\dots,tx^n)$ 作用于 $\mathbb{R}^n$ 不真：原点迷向群为 $\mathbb{R}^+$ 不紧；其他点轨道为开射线不闭。
\end{example}

\begin{theorem}{[商流形定理]}\label{theorem:ch21-商流形定理}
设 $M$ 是光滑流形，$G$ 在 $M$ 上光滑、自由、真作用。则 $M/G$ 是拓扑流形，并具有唯一光滑结构使 $\pi\colon M\to M/G$ 是光滑淹没。
\end{theorem}

\begin{proof}
令 $k=\dim G$，$m=\dim M$，$n=m-k$。唯一性：若两种光滑结构都使 $\pi$ 为淹没，由定理 4.29 恒等映射双向光滑，故相同。称坐标图 $(U,\varphi)$ \emph{适配}于 $G$-作用，若坐标 $(x,y)=(x^1,\dots,x^k,y^1,\dots,y^n)$ 使每条轨道与 $U$ 的交为空或单个 $y=c$ 切片。核心断言：每点有以它为中心的适配图。轨道是 $M$ 上分布 $D=\{X_p\}$（基本向量场）的积分流形；$D$ 光滑、对合，由 Frobenius 定理局部为切片。若不存在与每轨道至多交一切片的立方邻域 $U_0$，取缩小立方 $U_i$ 与横向 $n$-子流形 $Y$，得 $p_i,q_i\in Y_i$ 同轨道、$g_i p_i=q_i$，两序列收敛于 $p$；由真性取子列 $g_i\to g$，连续性给 $gp=p$，自由性给 $g=e$。映射 $\Psi\colon G\times Y\to M$ 在 $(e,p)$ 处微分为同构（$T_pM=T_p(Gp)\oplus T_pY$），故局部单射，与 $\Psi(g_i,p_i)=q_i=\Psi(e,q_i)$ 矛盾。故断言成立。$M/G$ Hausdorff（命题\cref{proposition:ch21-真作用Hausdorff}），适配图投影给出可数基与局部坐标，故为拓扑 $n$-流形；其坐标图册使 $\pi$ 的坐标表示为 $(x,y)\mapsto y$，是光滑淹没，过渡映射 $y^{new}=B(y)$ 光滑。唯一性由定理 4.29。
\end{proof}

\begin{figure}[ht]
\centering
\includegraphics[width=0.8\textwidth]{fig-21-1.png}
\caption{$M/G$ 的坐标图}\label{figure:ch21-quotient-chart}
\end{figure}
\end{document}
